\subsection{Pruebas unitarias y aceptación}

Como se mencionó en la sección de \nameref{sec:devops} y \nameref{sec:ci-cd}, se cuenta con un sistema de pruebas unitarias y de 
aceptación, para garantizar la calidad del código, y el correcto funcionamiento de las funcionalidades y servicios.
Utilizando \textbf{Go} para las pruebas unitarias, y \textbf{Behave} con \textbf{Python} para las pruebas de aceptación.
\vspace{0.25cm}

Las pruebas unitarias se encargan de probar cada uno de los módulos del \textit{core-service}, como lo son,
los \textit{services}, \textit{repositories}, \textit{utils} y \textit{middleware}.
\vspace{0.25cm}

Los repositorios se prueban utilizando una base de datos \textit{PostgreSQL} de prueba, que se levanta a través de \textit{Docker},
mientras que los servicios se prueban utilizando \textit{mocks} para simular el comportamiento de los repositorios, y así garantizar
que el servicio se comporte de la manera esperada, sin importar el comportamiento de los repositorios.
\vspace{0.25cm}

Para medir el porcentaje de cobertura de las pruebas unitarias, se utilizo la herramienta \textbf{CodeCov}, que permite 
generar un reporte de cobertura de código, y así garantizar que se están probando todas las partes del código. Dio como
resultado un porcentaje de cobertura del 75.95\%, lo que garantiza una buena calidad del código.
\vspace{0.25cm}

Por otro lado, las pruebas de aceptación se encargan de probar el correcto funcionamiento del servicio, a través de escenarios
que simulan el comportamiento de los usuarios, y así garantizar que el servicio se comporte de la manera esperada.
Para esto se utilizo, como se menciona, la herramienta \textbf{Behave} con \textbf{Python}.
\vspace{0.25cm}

Las pruebas de aceptación testean principalmente los endpoints del \textit{core-service}, para garantizar que se comporten 
correctamente, y que se cumplan los requisitos funcionales del servicio. La cantidad de escenarios de pruebas de aceptación, 
es de 54, lo que garantiza una buena cobertura de los requisitos funcionales del servicio.
\vspace{0.25cm}
